\chapter{Control de costos y ciclo de vida}

\section{Principios adoptados}

El prototipo se diseña para existir durante minutos u horas, no durante todo el curso. El costo se controla con cinco decisiones: cómputo optativo, una sola dirección pública, ausencia de NAT Gateway, volúmenes mínimos y destrucción automatizada. Las subredes, tablas de rutas, grupos de seguridad e Internet Gateway no justifican por sí solos mantener instancias encendidas.

AWS publica un cargo horario por cada dirección IPv4 pública, tanto asociada como ociosa \cite{aws-vpc-pricing}. Por ello el bastión se crea sólo al habilitar la prueba y se destruye después de recopilar la salida. El diseño tampoco reserva una Elastic IP, evitando que una dirección permanezca en la cuenta tras apagar una instancia.

\section{Componentes potencialmente facturables}

\begin{table}[H]
\centering
\caption{Superficie de costo del prototipo}
\begin{tabularx}{\textwidth}{p{4.0cm}Y p{3.0cm}}
\toprule
\textbf{Componente} & \textbf{Control} & \textbf{Duración esperada} \\
\midrule
Tres instancias EC2 pequeñas & Activación explícita y destrucción tras pruebas. & Menos de 2 horas \\
Tres volúmenes GP3 de 8 GiB & Eliminación junto con las instancias. & Menos de 2 horas \\
Una IPv4 pública & Sólo en el bastión; no se reserva Elastic IP. & Menos de 2 horas \\
Transferencia de datos & Únicamente administración y paquetes de prueba. & Volumen mínimo \\
VPC y subredes & No requieren cómputo activo; se eliminan al final. & Ventana de prueba \\
\bottomrule
\end{tabularx}
\end{table}

No se fija un precio total en el informe porque las tarifas, impuestos, créditos, región y condiciones de la cuenta pueden cambiar. Antes del \code{apply}, el equipo debe consultar la página vigente de precios y el panel de facturación. Cualquier crédito académico se trata como beneficio eventual, no como garantía.

\section{Etiquetado}

Todos los recursos incorporan:

\begin{itemize}
  \item \code{Project=CC3067-Fase-1}
  \item \code{Group=Grupo-1}
  \item \code{Environment=academic-poc}
  \item \code{ManagedBy=Terraform}
  \item \code{DeleteAfter=2026-10-19}
\end{itemize}

Las etiquetas permiten filtrar componentes y detectar residuos. AWS recomienda rastrear recursos durante todo su ciclo de vida y usar etiquetas para identificar activos que pueden retirarse \cite{aws-track-resources}.

\section{Lista de cierre}

Después de \code{terraform destroy} se revisa:

\begin{enumerate}
  \item cero instancias con etiqueta del proyecto;
  \item cero volúmenes EBS residuales;
  \item cero direcciones públicas o Elastic IP del proyecto;
  \item cero interfaces de red huérfanas;
  \item Internet Gateway separado y eliminado;
  \item VPC y subredes eliminadas;
  \item estado local protegido y no agregado a Git;
  \item archivo \code{destroy.txt} con el resultado final.
\end{enumerate}

\begin{warningbox}[Regla de operación]
Detener una instancia no equivale a destruir todos sus recursos. La puerta de salida exige revisar volúmenes, IP, interfaces y gateways, no solamente el estado de EC2.
\end{warningbox}
